% Analytically controlled family and an exact accessible-gap reduction.

\section{Parent path and accessible gap}
\label{sec:analytic-family}

The distinction between a global gap and a dynamically accessible gap is not
only a finite-instance effect.  We now give an even-cycle family for which
the symmetry reduction, the solution manifold, and an entire
parent-Hamiltonian path can be controlled.  The purpose of the family is a
separation theorem, not a claim that vertex cover on a cycle is classically
difficult.

\subsection{Even-cycle set cover and its joint symmetry}

Let the bases be the vertices of the even cycle $C_n$, $n=2k$, and let the
tasks be its edges.  Base $b$ covers the two edges incident on $b$.  A
$k$-subset $S\subset\mathbb Z_n$ is therefore a cover precisely when it is
one of the two alternating subsets.  The two solutions are exchanged by a
one-site rotation and form one orbit of the incidence automorphism group
$D_n$.

We work in the injective, register-symmetric subspace, with basis
$\{\ket S:|S|=k\}$.  Equivalently, this is the symmetric Dicke embedding of
the hard-core sector of the original registers.  Define
\begin{equation}
 U(S)=\#\{i\in\mathbb Z_n:i\notin S,\ i+1\notin S\}.
 \label{eq:cycle-energy}
\end{equation}
At half filling, $U(S)$ is also the number of adjacent occupied pairs.  Thus
$U=0$ exactly on the two minimum covers.

For $a\geq0$, consider the continuous-time Metropolis generator with a
Johnson-graph proposal,
\begin{equation}
 q_a(S,T)=\frac1n
 \min\{1,n^{-a[U(T)-U(S)]}\},
 \qquad |S\mathbin\triangle T|=2,
 \label{eq:johnson-metropolis}
\end{equation}
and $q_a(S,T)=0$ otherwise.  Its reversible measure is
\begin{equation}
 \pi_a(S)=Z_a^{-1}n^{-aU(S)}.
 \label{eq:cycle-gibbs}
\end{equation}
The symmetric-discriminant construction is standard for reversible Markov
chains and Gibbs-amplitude ground states
\cite{Szegedy2004,SommaBoixoBarnumKnill2008}.  Here $H_a$ has diagonal
$\sum_Tq_a(S,T)$ and off-diagonal entries
\begin{align}
 \langle T|H_a|S\rangle
 &=-\sqrt{q_a(S,T)q_a(T,S)}\\
 &=-\frac1n n^{-a|U(T)-U(S)|/2}.
 \label{eq:cycle-parent}
\end{align}
It is frustration free and has the known ground state
\begin{equation}
 \ket{\psi_a}=Z_a^{-1/2}\sum_{|S|=k}n^{-aU(S)/2}\ket S.
 \label{eq:cycle-parent-ground}
\end{equation}
At $a=0$, this is the Dicke state
$\ket{J_{n,k}}=\binom nk^{-1/2}\sum_{|S|=k}\ket S$, and the parent
Hamiltonian $H_{a=0}$ is exactly the Johnson-graph Laplacian at rate $1/n$.
It should not be confused with the restriction of $H_{\rm init}$ from
Eq.~\eqref{eq:hinit-rewrite}, which agrees with a Johnson Laplacian only up
to an additive constant after projection to the injective subspace.  The
whole parent path commutes with $D_n$.

This protocol has its own generator algebra and cyclic space, distinct from
Eq.~\eqref{eq:ku-rewrite}:
\begin{align}
 \mathfrak A_{\rm par}
 &=\operatorname{alg}^{*}\!\left(\{H_a:0\leq a\leq7\}\cup\{I\}\right),\\
 \mathcal K_J^{\rm par}
 &=\mathfrak A_{\rm par}\ket{J_{n,k}}.
 \label{eq:parent-cyclic-space}
\end{align}
Let $\mathcal H^{D_n}$ denote the $D_n$-fixed subspace of the half-filled
hard-core space and define
\begin{equation}
 \Delta_{\rm par}^{\rm cyc}(a)
 =\operatorname{gap}\!\left(H_a|_{\mathcal K_J^{\rm par}}\right).
 \label{eq:parent-cyclic-gap}
\end{equation}

\begin{proposition}[Cyclic accessibility of the parent path]
\label{prop:parent-cyclic-accessibility}
For every even $n=2k\geq6$ and $0\leq a\leq7$,
\begin{enumerate}
 \item $\mathcal K_J^{\rm par}$ is a common reducing subspace for the
 family $\{H_a\}_{0\leq a\leq7}$ and
 $\mathcal K_J^{\rm par}\subseteq\mathcal H^{D_n}$;
 \item the unique ground state $\ket{\psi_a}$ belongs to
 $\mathcal K_J^{\rm par}$; and
 \item
 \begin{equation}
  \Delta_{\rm par}^{\rm cyc}(a)\geq\Delta_{D_n}(a).
  \label{eq:parent-cyclic-fixed-comparison}
 \end{equation}
\end{enumerate}
Thus a lower bound proved in the $D_n$-fixed space is also a lower-bound
certificate for the dynamically accessible gap of this parent protocol.
\end{proposition}

Physically, the protocol explores no more than the $D_n$-fixed sector, and
restricting further to its true cyclic space cannot introduce an excitation
below the fixed-sector gap when the unique ground state is shared.  The full
argument is in Appendix~\ref{app:parent-accessibility}.

The number of configurations at energy $u$ is
\begin{equation}
 N_u=\frac{n}{k-u}\binom{k-1}{u}^2,
 \qquad 0\leq u<k.
 \label{eq:cycle-shell-count}
\end{equation}
At $a_f=7$, this shell count implies
\begin{equation}
 1-\langle\psi_{a_f}|\Pi_{U=0}|\psi_{a_f}\rangle
 =O(n^{-5}).
 \label{eq:cycle-endpoint-concentration}
\end{equation}
Appendix~\ref{app:parent-concentration} gives the run decomposition and
ratio bound.

\subsection{A uniform accessible-gap theorem}

We first separate the cycle geometry from the zero-range kinetics.

\begin{proposition}[Zero-range reduction]
\label{prop:cycle-zrp-reduction}
For every even $n=2k\geq6$ and every $a\geq0$, the gap of $H_a$ in the
$D_n$-fixed space obeys
\begin{equation}
 \Delta_{D_n}(a)
 \geq \frac{2}{n}[1-\cos(2\pi/k)]
 \operatorname{gap}\operatorname{ZRP}(K_k,r_a,k),
 \label{eq:cycle-zrp-reduction}
\end{equation}
where $r_a(0)=0$, $r_a(1)=n^{-a}$, $r_a(j)=1$ for $j\geq2$, and $K_k$
is the uniform mean-field kernel on the $k$ boxes.  Consequently, any
uniform polynomial lower bound for this mean-field gap gives a uniform
polynomial accessible-gap certificate for the parent path.
\end{proposition}

The factors in Eq.~\eqref{eq:cycle-zrp-reduction} have distinct physical
origins: $2/n$ is the update-rate normalization, $1-\cos(2\pi/k)$ is the
long-wavelength transport cost of the cycle, and the zero-range gap measures
the remaining occupation-redistribution bottleneck.  The exact isometry,
rate comparison, and normalization are proved in
Appendix~\ref{app:normalized-gap-proof}.

\begin{proposition}[Uniform accessible gap]
\label{prop:cycle-uniform-gap}
For every even $n=2k\geq6$ and every $a\geq0$,
\begin{align}
 \operatorname{gap}\operatorname{ZRP}(K_k,r_a,k)
 &\geq 2n^{-a}k^{-3},
 \label{eq:mean-field-bound}\\
 \Delta_{D_n}(a)&\geq1024\,n^{-(a+6)}.
 \label{eq:uniform-accessible-gap}
\end{align}
Consequently, along $0\leq a\leq7$,
$\min_a\Delta_{D_n}(a)\geq1024\,n^{-13}$.
\end{proposition}

Physically, the polynomial bound accumulates three conservative costs:
normalized subset updates, long-wavelength transport around the cycle, and
redistribution of vacancies between effective boxes.  Its exponent is a
certificate, not a prediction of the finite-size gap scaling.  The complete
comparison is in Appendix~\ref{app:normalized-gap-proof}.

The gap theorem can be converted into an adiabatic statement in the
abstract Hamiltonian-evolution model.  This does not by itself supply a
gate-level implementation of the path.

\begin{corollary}[Conditional adiabatic preparation]
\label{cor:cycle-adiabatic-preparation}
In the abstract Hamiltonian-access model, let $a(s)=7s$, $0\leq s\leq1$,
and evolve from $\ket{J_{n,k}}$ under the restricted path
$H_{a(s)}|_{\mathcal K_J^{\rm par}}$ for total time $T$.  There is a
universal constant $C>0$ such that, for every $0<\epsilon<1$, the sufficient
condition
\begin{equation}
 T\geq C\epsilon^{-1}n^{41}(\log n)^2
 \label{eq:cycle-adiabatic-runtime}
\end{equation}
produces a final state at distance $O(\epsilon)$ from
$\ket{\psi_7}$, up to phase.  Its cover probability is therefore
$1-O(n^{-5})-O(\epsilon)$.  This runtime excludes the costs of preparing the
Dicke state and realizing or simulating $H_{a(s)}$.
\end{corollary}

Physically, the polynomial cyclic gap prevents dark global crossings from
setting the preparation time of this protocol.  The large runtime exponent
instead reflects deliberately conservative derivative and inverse-gap
bounds and does not include Hamiltonian-simulation or state-preparation
costs.  The complete adiabatic and measurement estimates are in
Appendix~\ref{app:parent-adiabatic}.

Exact $D_n$-orbit quotients for $n=6,8,\ldots,18$ show finite-size gaps much
larger than the conservative theorem, but are not used to infer scaling.
The grid and all intermediate values are reported in
Appendix~\ref{app:numerical-protocol} and the archived data.

The path also clarifies what is and is not inherited from the original
multi-register Hamiltonian.  Its endpoint ground state has the same
solution-producing measurement objective, but the initial ground state is
the injective Dicke state rather than the product-register uniform state,
and the endpoint is the frustration-free Gibbs parent $H_{a_f}$ rather than
the original $H_{\mathrm{prob}}$.  No interpolation between these two
endpoints is assumed.  The exact dark multiplicity closure of the original
linear interpolation and the uniformly polynomially gapped parent path are
therefore complementary statements
about the same symmetric instance family, not two descriptions of one
protocol.
